Apply (O.12) at coefficient one.  Put $V=U\setminus\{p\}$,
$v=\prod_{r\in V}r$, so $n_U=n_V+1$.  For a fresh prime $r$ set
\[
 C_p(vr)=\frac{a_p-(\psi/\eta)(p)(-vr/p)-1}{2}\in\Z,
 \qquad c_p(vr)=C_p(vr)\bmod2.
\]
The two surviving operator terms at $pr$ and the good recurrence give
\[
 H_U[pr]\equiv
   \eta(p)C_p(vr)\,2^{-n_V-1}a(vr)\pmod{2\mathcal O_{\rm coef}}.
\]
For the even source this is an equality before reduction.  For the
odd source the error is in $2^{M_\theta-n_U-1}\mathcal O_{\rm coef}$,
which is contained in $2\mathcal O_{\rm coef}$ at our chosen stage.
Choose fresh $q,q_\tau,r_0,r_1$ representing respectively
$g,g\tau,g\phi,g\tau\phi$ in one sufficiently precise joint quotient
for the $U$ and $V$ forms and all retained characters.
The omitted form uses the same $g,\eta,\psi,\delta$ and stage,
and its level omits the entire rational prime $p$.
Thus the two $V$ symbols at $vr_0,vr_1$ agree.  Their Legendre
signs at $p$ are opposite, so $c_p(vr_0)+c_p(vr_1)=1$ in $\F_2$.
The coefficient-one reading and the reduced trace identity now give
\[
             \mathcal A_U(Uq_\tau)+\mathcal A_U(Uq)=\mathcal A_V(Vr_0)\qquad\text{in }\F_2.
\]
Here the subscript records the operator set, not a new coefficient
source or minimum. The factor $\eta(p)$ reduces to one; no factor
two is cancelled in (O.12). For every retained quadratic character
$\chi$ unramified at $p$, the representative of the omitted vertex
satisfies
\[
 \chi(g\phi)=\chi(g)\chi(\phi)
            =\chi(\operatorname{Fr}_q)\chi(\operatorname{Fr}_p).
\]
For a fixed-support character this is the product of the two vertex
labels. For the character of a third vertex $r\in V$, it adds the
two bits $(r/q)$ and $(r/p)$ in $\F_2$. Reciprocity then gives the
reverse bits from the same endpoint types. All additional stationary
tests were included in the joint quotient, so they obey this formula
as well. This proves the full product-label omission rule, including
every external incidence.

A fresh vertex with trivial label in the elementary extension for
$S\cup U$ has the same reduced reading as $t(1)$, hence gives zero.
An isolated primary vertex has precisely this fixed label and these
zero incidence bits. A contracted cluster of primary vertices has
the same property when its external incidences are neutral. The
result is still a nonempty configuration, including the one-vertex
case; no contraction deletes the last vertex.

